\documentclass[11pt,a4paper]{article}
\usepackage[T1]{fontenc}
\usepackage{isabelle,isabellesym}
\usepackage{amsmath}
\usepackage{amssymb}
\usepackage{url}

\isabellestyle{it}

\begin{document}

\title{A Minimal Homotopy Type Theory Object Logic for Isabelle/Pure}
\author{Arthur Freitas Ramos \and David Barros Hulak \and
  Ruy J. G. B. de Queiroz}
\date{July 24, 2026}
\maketitle

\begin{abstract}
This entry presents an independently authored higher-order abstract syntax
encoding of intensional Martin-L\"of type theory
\cite{MartinLof1984,NordstromPeterssonSmith1990} and a
core homotopy type theory library \cite{HoTTBook2013} in Isabelle/Pure. The
primitive object-logic rules are explicit axioms. Derived results cover
identity-path algebra, transport, homotopies, contractibility, fibers,
bi-invertibility, and the identity, inverse, and composite equivalences, so
type equivalence is proved reflexive, symmetric, and transitive. The main
characterization proves that a map is bi-invertible if and only if all of
its fibers are contractible; contractibility is also shown to transfer along an
equivalence in both directions.
Function extensionality and univalence are isolated as optional additional
axioms, each stated as bi-invertibility of its canonical map.
\end{abstract}

\tableofcontents

\section{Scope and independence}

The source was written from the mathematical presentation in the Homotopy Type
Theory book \cite{HoTTBook2013}. It does not import, copy, rename, or adapt
Joshua Chen's Isabelle/HoTT implementation. Chen's earlier Isabelle
formalization \cite{Chen2021} is cited as related work, not as a source
dependency.

The distinct archival value is a minimal, automation-free, explicitly
axiomatized HoTT object logic whose complete trusted boundary is visible in
ordinary Isabelle theory files. It is intended as a reference foundation and
audit target rather than as a replacement for a feature-rich interactive HoTT
environment. That role is complementary to the large univalent libraries built
on proof assistants with native dependent types, such as the Coq HoTT library
\cite{BauerEtAl2017} and UniMath \cite{Voevodsky2015}, which optimize for
mathematical reach rather than for a small auditable rule set.

\begin{center}
\small
\begin{tabular}{@{}p{0.20\textwidth}p{0.36\textwidth}p{0.36\textwidth}@{}}
\hline
Aspect & This entry & Chen's Isabelle/HoTT \cite{Chen2021} \\
\hline
Development &
Independently authored object logic from published mathematical rules &
The original Isabelle/HoTT implementation \\
Implementation &
Ordinary Isabelle/Pure theories; no custom Isabelle/ML automation &
Bidirectional type checking and elaboration infrastructure \\
Trust presentation &
Primitive rules and optional extensional axioms are visibly separated &
Richer implementation and proof infrastructure for interactive use \\
Formal scope &
Focused foundational library: the full contractible-fiber characterization of
equivalences, and closure under identity, inversion, and composition &
Broader foundational results and a more convenient working environment \\
Intended use &
Explicit reference foundation and audit surface &
Interactive HoTT formalization and elaboration environment \\
\hline
\end{tabular}
\end{center}

\section{Representation}

Object expressions inhabit a fresh Pure type. Object binders are represented
with higher-order abstract syntax \cite{PfenningElliott1988}, so Pure beta
substitution implements object-level substitution. Isabelle/Pure is used as a
logical framework in the sense of Paulson \cite{Paulson1989,%
WenzelPaulsonNipkow2008}; the presentation
also follows the general logical-framework discipline of
Harper, Honsell, and Plotkin \cite{HarperHonsellPlotkin1993}. The object logic
has a dedicated typing judgment and a primitive type of universe levels. A
zero level, a successor operation, and a strict universe-level relation
index cumulative Russell-style universes. The relation is axiomatized by a
successor step, transitivity, and irreflexivity. Pure has no falsity, so
irreflexivity is stated in ex-falso form: a self-relation \(i \prec_l i\)
derives an arbitrary Pure proposition. Asymmetry, and hence the absence of
short successor cycles, is derived from those axioms. The intended model reads
levels as natural numbers, so the strictness axiom is not vacuous. The entry
does not identify levels with Isabelle natural numbers or postulate an
exhaustion principle for them.

Dependent products and sums use Pure functions for their families. Application
and projections are object-language constants, while their beta equations are
Pure equalities. The empty type has formation and elimination but no
introduction rule. Unit and natural-number eliminators have the expected
computation equations. Identity types have reflexivity and based path
induction.

\section{Derived identity and equivalence layer}

Transport is defined by path induction. Path inversion is another direct
instance of path induction, and composition transports the first path along
the second in a family of paths with fixed source. The development constructs
inhabitants of the left and right unit laws, associativity, and both inverse
laws, which are the groupoid laws of the identity type in the sense of
Hofmann and Streicher \cite{HofmannStreicher1998}. Functorial action on paths
is obtained by induction, and homotopies are dependent products of pointwise
identity types.

A contractible type packages a center and a path from that center to every
point. Fibers are dependent pairs of preimages and image paths, and a
contractible map has a contraction in every fiber. Left and right inverse data
are kept independent; their product is bi-invertibility. An equivalence is a
function equipped with such a witness.

The development proves both directions of the standard characterization. From
contractible fibers it selects each fiber center as an inverse image; applying
the first projection to the fiber contraction yields the left-inverse
homotopy. Conversely, independent left and right inverses are combined into a
single quasi-inverse and upgraded to coherent half-adjoint data. The triangle
identity then supplies the dependent second-component path needed to contract
each fiber. This proof derives functorial laws for \texttt{ap}, path
whiskering and cancellation, homotopy naturality, and paths of dependent pairs.
No new primitive rule is introduced for the characterization.

The identity function is proved bi-invertible using reflexive homotopies. For
composition, the selected left inverses and selected right inverses are
composed separately. Their new inverse homotopies apply the outer inverse or
map to an existing pointwise path and then concatenate with the remaining
inverse path. Thus closure of equivalences under composition is derived rather
than postulated.

Inversion is derived from the same quasi-inverse construction. Reading the
retraction and section homotopies of a bi-invertible \(f\) from the other side
exhibits \(f\) itself as both a left and a right inverse of the selected
inverse \(g\), so \(g\) is bi-invertible. A general constructor turns any
quasi-inverse triple into a bi-invertibility witness, which is the counterpart
of the projection lemmas that read the inverse map and inverse homotopy out of
an existing witness. Identity, inversion, and composition together make type
equivalence reflexive, symmetric, and transitive, each witnessed by an explicit
object-language term rather than by a meta-level assertion. Combining inversion
with the preservation result also transfers contractibility backwards along an
equivalence, so contractibility is invariant in both directions.

\section{Extensional principles}

For two dependent functions, \texttt{happly} sends a path of functions to the
pointwise homotopy by path induction. Function extensionality is the additional
axiom that this canonical map is bi-invertible. The exported \texttt{funext}
term is selected from the left inverse supplied by that witness.

For types \(A,B\) in a universe, \texttt{idtoequiv} sends a universe path to an
equivalence by path induction, with the identity equivalence as its reflexivity
case. Univalence is the additional axiom that this canonical map is
bi-invertible. The function-extensionality and univalence theories both import
the intensional equivalence layer directly; neither optional axiom is hidden in
the earlier theories.

\section{Trust boundary}

The typing judgment, universe hierarchy, primitive type formers, eliminators,
and their computation equations are axiomatized in Isabelle/Pure. Isabelle
checks every derived theorem relative to those axioms. Function extensionality
and univalence extend that boundary only in their respective final theories.
There is no custom Isabelle/ML proof automation and no oracle. Irreflexivity of
the universe-level relation is encoded in Pure's ex-falso form: a hypothetical
self-relation derives an arbitrary Pure proposition. That premise is impossible
in the intended model, which reads levels as natural numbers, \(\prec_l\) as
the strict order, and the successor as \(+1\); every level axiom holds there,
no self-relation does, and the axiom is used only to derive level asymmetry.

The entry does not prove normalization, canonicity, consistency, decidability
of type checking, or a semantic model such as the simplicial
\cite{KapulkinLumsdaine2021} or cubical \cite{BezemCoquandHuber2014} models of
univalent foundations. It also does not include higher inductive types, cubical
computation, a general inductive-definition package, or automatic universe
inference.

\section{Theory structure}

{\sloppy
\begin{itemize}
\item \texttt{HoTT\_Base} introduces expressions, typing, and cumulative
  universes.
\item \texttt{HoTT\_Type\_Formers} introduces dependent products and sums,
  empty and unit types, and natural numbers.
\item \texttt{HoTT\_Identity} derives transport and path algebra from path
  induction.
\item \texttt{HoTT\_Equivalence} develops contractibility, bi-invertibility,
  the identity equivalence, and composition of equivalences.
\item \texttt{HoTT\_Equivalence\_Characterization} proves the bidirectional
  contractible-fiber criterion and that equivalences preserve contractibility.
\item \texttt{HoTT\_Equivalence\_Inverse} inverts equivalences and collects the
  reflexivity, symmetry, and transitivity witnesses.
\item \texttt{HoTT\_Function\_Extensionality} adds function extensionality as
  an optional axiom.
\item \texttt{HoTT\_Univalence} adds univalence as an optional axiom.
\end{itemize}
}

Both extensional principles live in separate final theories, so their
additional assumptions remain visible.

\section{Artifact and reproducibility}

The artifact targets Isabelle2025-2 and the session
\texttt{HoTT\_Object\_Logic}. It contains eight top-level theory files and no
Isabelle/ML source. A proof-only build is:

{\small
\begin{verbatim}
.\tools\build.ps1 -Project hott-object-logic-isabelle -NoDocument
\end{verbatim}
}

A clean proof and document build, followed by the AFP-style source audit, is:

{\small
\begin{verbatim}
.\tools\build.ps1 -Project hott-object-logic-isabelle -Clean
.\tools\audit.ps1 -Project hott-object-logic-isabelle
\end{verbatim}
}

The equivalent portable invocation is
\texttt{isabelle build -c -D <entry directory> HoTT\_Object\_Logic}.

The final audit records the exact source-line and theorem counts and checks for
\texttt{sorry}, \texttt{oops}, missing AFP files, and generated artifacts. For
this revision it reports eight theories, 132 object-logic statements, twelve
headline theorems, 7,609 theory-source lines, no forbidden proof tokens, and no
findings. Every proof is structured Isar: the entry contains no \texttt{apply}
script, no \texttt{back} or \texttt{apply\_end}, no diagnostic or
counterexample command, no \texttt{smt} invocation, no unnamed statement
carrying a global attribute, and no use of quick-and-dirty mode. A clean
AFP-style proof and PDF build takes well under a minute on a contemporary
laptop; build times are machine-dependent.

\input{session}

\bibliographystyle{abbrv}
\bibliography{root}

\end{document}
